\documentclass[10pt,a4paper]{amsart}
\usepackage[margin=19mm]{geometry}
\usepackage{amsmath,amssymb,mathtools}
\usepackage[scr=rsfs,cal=euler]{mathalfa}
\usepackage{xfrac}
\usepackage[unicode,hidelinks,bookmarks=false]{hyperref}
\newcommand{\ie}{i.e.\ }
\newcommand{\cf}{cf.\ }
\newcommand{\resp}{respectively\ }
\newcommand{\Subsection}{\subsection}
\providecommand{\nobreakdash}{\nobreak\hbox{-}}
\setlength{\emergencystretch}{2em}
\multlinegap0pt
\usepackage{bm}
\usepackage{bbm}
\usepackage{dsfont}
\usepackage{xspace}
\usepackage{cancel}
\usepackage{mathtools}
\usepackage{amsthm}
\makeatletter
\@ifundefined{theo}{\newtheorem{theo}{Theo}}{}
\@ifundefined{prop}{\newtheorem{prop}[theo]{Proposition}}{}
\makeatother
\newcommand{\bU}{\mathbb{U}}
\newcommand{\bW}{\mathbb{W}}
\title[Cluster parking functions II: $q,t$-dihedral sieving via dia]{Cluster parking functions II: $q,t$-dihedral sieving via diagonal coinvariants}
\author{Matthieu Josuat-Verg\`es}
\date{Source: arXiv:2607.04999v1 (CC BY 4.0)}
\begin{document}
\maketitle

\noindent\emph{Source and licence.} Cluster parking functions II: $q,t$-dihedral sieving via diagonal coinvariants. Version arXiv:2607.04999v1 (CC BY 4.0), distributed under the Creative Commons Attribution 4.0 International licence, as recorded in the arXiv licence metadata for this exact version and shown on the abstract page https://arxiv.org/abs/2607.04999v1. Full rights notice: licenses/arxiv-2607.04999-CC-BY-4.0.txt.\par\smallskip

\noindent\textit{Editorial scope.} Counted unit: the Prop whose source label is \texttt{prop:algparkessai} in the article (source0.tex lines 1681--1688, source label \texttt{prop:algparkessai}), delivered together with the article's own complete original proof of it (source0.tex lines 1690--1753, one \texttt{proof} environment of 64 lines). No mathematics is written, repaired or re-derived: statement and proof are the article's own text line for line. Required context retained uncounted: none. Every definition, display and label of those lines is retained unchanged. Omitted from this package and not redistributed: 1--1680, 1689--1689, 1754--1937. Nothing inside a retained excerpt is omitted or altered: every retained body is delivered exactly as the article's lines are written, so no omission receipt of any kind accompanies this package. Citations: the retained excerpts carry no citation command at all, so no bibliography entry of the article is transcribed and no reference is redistributed. Reference adaptation: none was needed and none was made; every reference inside the delivered unit resolves inside it, so no unresolved reference or citation is produced. Printed slips kept literal: no printed word, symbol, hypothesis or number of the counted statement or of its proof was altered. Layout and class: the article's own class is \texttt{article} with packages including babel, authblk, amsmath, amssymb, amsthm, mathtools, hyperref, tikz, physics, a4wide, caption, enumitem, url; the wrapper is the lane's pinned \texttt{amsart} editorial wrapper at 10pt on A4, which restates the theorem-like environments the retained text needs and adds the article's own macro definitions that the retained text needs (\texttt{\textbackslash bU}, \texttt{\textbackslash bW}), with extra wrapper packages (bm, bbm, dsfont, xspace, cancel, mathtools). The delivered numbering is the unit's own. Assets: no retained span contains a figure environment, a graphics-inclusion command, a PDF-page inclusion, a table or a TikZ picture; no image file is copied into this package, the package declares no asset dependency and no image file is redistributed with it. Licence: version arXiv:2607.04999v1 of this article is distributed under the Creative Commons Attribution 4.0 International licence, as recorded in the arXiv licence metadata for this exact version and shown on the abstract page https://arxiv.org/abs/2607.04999v1. A full rights notice is delivered at licenses/arxiv-2607.04999-CC-BY-4.0.txt. Characters: every retained excerpt is pure ASCII, so the delivered engine, XeLaTeX with its default Latin Modern text font, reports no missing character.
\begin{prop} \label{prop:algparkessai}
    The character of $\bU_{k+2} \times \mathfrak{S}_n$ acting on $\bW_k^{\otimes n}$ is
    \[
        (\mu, \sigma)
            \mapsto
        (-1)^n \cdot \epsilon(\sigma) \cdot (-k-1)^{\dim\ker(\sigma-\mu I)}.
    \]
\end{prop}

\begin{proof}
A basis of $\bW_k$ is given by the functions $f(z) = z^i$ for $1\leq i \leq k+1$ (for example, by the finite Fourier transform).  So, a basis of the tensor power $\bW_k^{\otimes n}$ is given by the elements
\[
    z^{\otimes I} := z^{i_1} \otimes \dots \otimes z^{i_n},
\]
where $I = (i_1,\dots,i_n) \in \{1,\dots,k+1\}^n$ is a multi-index. 

The action of $\mu\in\bU_{k+2}$ on the basis $(z^i)_{1 \leq i \leq k+1}$ is diagonal, since
\[
    \mu \cdot z^i = (\mu z)^i = \mu^i z^i.
\]
It follows that the action of $\bU_{k+2}$ on the tensor power $\bW_k^{\otimes n}$ is also diagonal, namely:
\[
    \mu \cdot z^{\otimes I} 
        =
    (\mu^{\sum I}) z^{\otimes I}
\]
where $\sum I = \sum_{j=1}^n i_j$.  On the other side, the natural action of $\mathfrak{S}_n$ permutes the basis elements:
\[
    \sigma \cdot z^{\otimes I}
        =
    z^{\otimes (\sigma\cdot I)}
\]
where $\sigma \cdot I = (i_{\sigma^{-1}(1)}, \dots, i_{\sigma^{-1}(n)} )$.  All in all, we see that the action of $\bU_{k+2} \times \mathfrak{S}_n$ on $\bW_k^{\otimes n}$ is monomial:
\[
    (\mu,\sigma) \cdot z^{\otimes I}
        =
    (\mu^{\sum I}) z^{\otimes (\sigma\cdot I)}.
\]
We thus get the trace of $(\mu,\sigma)$ as a sum:
\[
    \operatorname{Tr}(\mu,\sigma)
        =
    \sum_{\substack{I \in \{1,\dots,k+1\}^n, \\ \sigma\cdot I = I}} \mu^{\sum I}.
\]
Let us represent the cyclic type of $\sigma$ by an integer partition $(\lambda_1,\dots,\lambda_\ell)$ of $n$.  The sum can be rewritten:
\begin{align*}
    \operatorname{Tr}(\mu,\sigma)
        &=
    \sum_{(i_1,\dots,i_\ell) \in \{1,\dots,k+1\}^\ell} \mu^{\lambda_1 i_1 + \dots + \lambda_\ell i_\ell} \\[2mm]
        &=
    \prod_{u=1}^\ell \bigg(
        \sum_{i_u = 1}^{k+1} \mu^{\lambda_u i_u}.
    \bigg)
\end{align*}
To evaluate the sum between parentheses, note that:
\begin{itemize}
    \item if $\mu^{\lambda_u} = 1$ ({\it i.e.}, the order of $\mu$ divides $\lambda_u$), the sum is $k+1$,
    \item otherwise, $\mu^{\lambda_u}$ is a $(k+2)$nd root of unity different from $1$, so that 
    \[
        \sum_{i_u = 0}^{k+1} \mu^{\lambda_u i_u}
            = 
        0,
    \]
    and the sum from $1$ to $k+1$ is $-1$.
\end{itemize}
So we get:
\begin{align*}
    \operatorname{Tr}(\mu,\sigma)
        &=
    (-1)^{\ell - \dim\ker(\sigma-\mu I)} (k+1)^{\dim\ker(\sigma-\mu I)}.
\end{align*}
It remains to notice that $\epsilon(\sigma) = (-1)^{n-\ell}$.
\end{proof}

\end{document}
